% Input-ready conventional-proof component.  Labels use the prefix sp:.
% Requires the mean_field.tex component and its mf: labels.

\section{A spectral obstruction with an explicit positive margin}
\label{sp:section}

The full score in Theorem~\ref{mf:criterion} admits an operator
factorization that gives an obstruction in every fixed latent dimension.
Throughout this section the rows and empirical measures satisfy the
bounded-row hypotheses of Section~\ref{mf:setup}. Eigenvalues of the
limiting covariance \(\Sigma\) are listed in decreasing order, with
missing eigenvalues interpreted as zero. All analytic arguments in
this section are conventional proofs.

\begin{theorem}[At most two supercritical covariance directions]
\label{sp:main}
Let \(r\ge1\), \(t>0\), and \(\|v_{i,N}\|\le L\), and suppose
\[
 \Sigma_N=\frac1N V_N^TV_N\longrightarrow\Sigma,\qquad
 t\lambda_1(\Sigma)>1,\qquad t\lambda_3(\Sigma)\le1.
\]
For \(A_N=I_N+(t/N)V_NV_N^T\), every sequence of row orders satisfies
\begin{equation}\label{sp:positive-margin}
 \liminf_{N\to\infty}
 \frac{D_1(A_N)}{N(N-1)\per A_N}
 \ge
 \frac{\bigl(t\lambda_1(\Sigma)-1\bigr)^2}{64r\,t^2L^4}>0.
\end{equation}
In particular the endpoint derivative is eventually positive.
The hypothesis includes the boundary \(t\lambda_3(\Sigma)=1\).
\end{theorem}

The nontrivial eigenvalues of \(A_N\) converge to
\(1+t\lambda_j(\Sigma)\). Thus the theorem permits at most two
limiting eigenvalues strictly above two, while requiring the largest
one to exceed two. This is a statement about the specified bounded-row,
fixed-rank mean-field families. Its conclusion does not assert a
spectral criterion for unrestricted finite matrices.

\subsection{The full score as a Hilbert--Schmidt norm}

Fix an ordered limit \(\mu\) as in \eqref{mf:empirical} and a point
\(h\in\mathbb C^r\). Use \(p=p_h\), \(q=q_h\), and \(m=m_h\)
from \eqref{mf:p}--\eqref{mf:q}, and define the real symmetric matrix
\begin{equation}\label{sp:covariance-resolvent}
 S_h=t\int\frac{vv^T}{1+t|v\cdot h|^2}\,d\mu(x,v).
\end{equation}
Let \(H=L^2(\mu;\mathbb C)\). On this space define
\[
 (\mathsf Kf)(x,v)=\int\operatorname{sgn}(y-x)f(y,w)\,d\mu(y,w),
 \qquad \mathsf Qf=\int qf\,d\mu\in\mathbb C^r,
\]
and set
\[
 \mathsf D=M_{\sqrt{1-p}},\qquad
 \mathsf T:H\longrightarrow\mathbb C^r\oplus H,
 \qquad \mathsf Tf=(\mathsf Qf,\mathsf Df).
\]
For an operator between these complex Hilbert spaces, write
\(\mathsf T^{\mathrm{tr}}=J_H\mathsf T^*J_{\mathbb C^r\oplus H}\),
where \(J\) is coordinatewise conjugation. This transpose has the same
operator norm as \(\mathsf T\).

\begin{lemma}[Exact factorization]\label{sp:factorization}
The sign-kernel operator satisfies \(\|\mathsf K\|_{\mathrm{HS}}^2=1\),
and the full score \eqref{mf:score} is
\begin{equation}\label{sp:factorization-eq}
 \mathcal R(h)=
 \|\mathsf T\mathsf K\mathsf T^{\mathrm{tr}}\|_{\mathrm{HS}}^2.
\end{equation}
If \(S_h\preceq I_r\), then \(\mathsf T^*\mathsf T\preceq I_H\).
\end{lemma}
\begin{proof}
The diagonal set \(x=y\) has \(\mu\otimes\mu\)-measure zero because
the position marginal is Lebesgue. Squaring the sign kernel therefore
gives Hilbert--Schmidt norm squared one.

The old--old block \(\mathsf Q\mathsf K\mathsf Q^{\mathrm{tr}}\)
has entries
\[
 \iint\operatorname{sgn}(y-x)q_a(x,v)q_b(y,w)\,d\mu(x,v)d\mu(y,w).
\]
It is the antisymmetric matrix of the ordered wedge \(W_h\). Its
matrix Hilbert--Schmidt norm squared is \(2\|W_h\|^2\).
Each of the two mixed blocks has squared norm
\[
 \int(1-p(v))\|2Q_h(x)-Q_h\|^2\,d\mu(x,v).
\]
The noise--noise block has kernel
\(\sqrt{1-p(v)}\operatorname{sgn}(y-x)\sqrt{1-p(w)}\), so its
squared norm is \((1-m)^2\). Adding the four block norms proves
\eqref{sp:factorization-eq}, including the factor two in the mixed term.

For the contraction assertion, put \(b=q/\sqrt p\) on \(\{p>0\}\)
and \(b=0\) elsewhere. Since \(q=0\) on \(\{p=0\}\),
\[
 \int bb^*\,d\mu
 =t\int_{\{p>0\}}(1-p)vv^T\,d\mu\preceq S_h.
\]
This is an inequality: rows with zero projection on \(h\) can still
contribute to \(S_h\). Weighted Cauchy--Schwarz consequently gives
\[
 \|\mathsf Qf\|^2
 =\left\|\int b\sqrt p\,f\,d\mu\right\|^2
 \le\|S_h\|\int p|f|^2\,d\mu
 \le\int p|f|^2\,d\mu.
\]
Adding \(\|\mathsf Df\|^2=\int(1-p)|f|^2\,d\mu\) proves the claim.
\end{proof}

\begin{lemma}[Quantitative contraction loss]\label{sp:gap}
If \(S_h\preceq I_r\), then
\begin{equation}\label{sp:gap-eq}
 \mathcal R(h)\le1-\frac{m_h^2}{16r}.
\end{equation}
\end{lemma}
\begin{proof}
Put
\(\mathsf P=I_H-\mathsf T^*\mathsf T=M_p-\mathsf Q^*\mathsf Q\succeq0\).
The Hilbert--Schmidt ideal inequality and
Lemma~\ref{sp:factorization} imply
\begin{equation}\label{sp:trace-loss}
 \mathcal R(h)
 \le\|\mathsf T\mathsf K\|_{\mathrm{HS}}^2
 =1-\operatorname{Tr}(\mathsf K^*\mathsf P\mathsf K).
\end{equation}
The conclusion follows immediately when \(m=0\), so assume \(m>0\).
Choose an integer \(d>r\) and write \(w=m/d\). The cumulative
\(p\)-mass is continuous, and \(0\le p\le1\). Hence \([0,1]\)
can be partitioned into \(d\) consecutive intervals of \(p\)-mass
exactly \(w\), each of length at least \(w\).

On interval \(j\), let \(g_j\) rise linearly from zero to one over
its first segment of length \(w/4\), remain one, and fall linearly to
zero over its last segment of length \(w/4\). Extend it by zero and
put \(f_j=-g'_j/2\). Values at the finitely many corners do not matter.
Integration and the Lebesgue position marginal give
\begin{equation}\label{sp:hats}
 \mathsf Kf_j=g_j,\qquad
 \|f_j\|^2=\frac2w,\qquad
 \int p|g_j|^2\,d\mu\ge\frac w2.
\end{equation}
For the last bound, discard the two edge strips. Their total length
is \(w/2\), so they contain at most \(w/2\) of \(p\)-mass; the
remaining mass lies on the plateau.

The disjoint supports make the \(f_j\) orthogonal. Their span \(E\)
has dimension \(d\), and \eqref{sp:hats} gives, for every \(f\in E\),
\[
 \int p|\mathsf Kf|^2\,d\mu\ge\frac{w^2}{4}\|f\|^2.
\]
The subspace \(E\cap\ker(\mathsf Q\mathsf K)\) has dimension at
least \(d-r\). Each unit vector in that subspace satisfies
\[
 \langle\mathsf K^*\mathsf P\mathsf Kf,f\rangle
 =\int p|\mathsf Kf|^2\,d\mu-\|\mathsf Q\mathsf Kf\|^2
 \ge\frac{w^2}{4}.
\]
Summing over \(d-r\) orthonormal vectors and using positivity of the
trace-class operator gives
\[
 \operatorname{Tr}(\mathsf K^*\mathsf P\mathsf K)
 \ge\frac{(d-r)m^2}{4d^2}.
\]
Take \(d=2r\) in \eqref{sp:trace-loss} to obtain
\eqref{sp:gap-eq}.
\end{proof}

\subsection{Stationarity and the third covariance eigenvalue}

\begin{lemma}[An inertia bound]\label{sp:inertia}
Suppose \(h\) is a stationary point of \(\Phi\) and
\(\operatorname{Re}h,\operatorname{Im}h\) are linearly independent.
If \(S_h\) has \(u\) eigenvalues strictly greater than one, counted
with multiplicity, then \(t\Sigma\) has at least \(u+2\) eigenvalues
strictly greater than one.
\end{lemma}
\begin{proof}
Differentiation gives \(S_hh=h\). Since \(S_h\) is real symmetric,
it fixes the two-dimensional real space
\(E=\operatorname{span}_{\mathbb R}(\operatorname{Re}h,
\operatorname{Im}h)\) pointwise. Let \(U\) be the sum of its
eigenspaces with eigenvalue greater than one. Then
\(\dim U=u\) and \(U\perp E\).

The covariance correction
\begin{equation}\label{sp:correction}
 C=t\Sigma-S_h=t\int p_h(v)vv^T\,d\mu\succeq0
\end{equation}
is strictly positive on every nonzero vector \(x\in E\). Indeed,
\(x^TS_hx=\|x\|^2>0\) implies positive measure for
\(\{v:v\cdot x\ne0\}\). On this set \(v\cdot h\ne0\), since
\(x\) is a real linear combination of the real and imaginary parts
of \(h\). Thus \(p_h(v)>0\) there and \(x^TCx>0\).

For a nonzero \(x\in E\oplus U\), a nonzero \(U\)-component gives
\(x^T(S_h-I)x>0\), while \(x^TCx\ge0\). If that component
vanishes, \(x\in E\setminus\{0\}\) and the strict positivity
comes from \(C\). Hence \(t\Sigma-I\) is positive definite on
the \((u+2)\)-dimensional space \(E\oplus U\). The min--max
principle proves the claimed strict eigenvalue count.
\end{proof}

\begin{lemma}\label{sp:maximizer-contraction}
If \(t\lambda_3(\Sigma)\le1\), then \(S_h\preceq I_r\) at
every global maximizer \(h\) of \(\Phi\).
\end{lemma}
\begin{proof}
If \(h\) is real up to a global phase, rotate that phase to make it
real. For every real vector \(z\), the imaginary variation gives
\[
 \frac12\left.\frac{d^2}{d\epsilon^2}
       \Phi(h+i\epsilon z)\right|_{\epsilon=0}
 =z^T(S_h-I)z\le0.
\]
Therefore \(S_h\preceq I\); this also covers \(h=0\).
Otherwise the real and imaginary parts are independent. An eigenvalue
of \(S_h\) above one would, by Lemma~\ref{sp:inertia}, force
\(t\lambda_3(\Sigma)>1\), a contradiction.
\end{proof}

\begin{proof}[Proof of Theorem~\ref{sp:main}]
First suppose an ordered empirical limit exists. At every maximizer,
Lemma~\ref{sp:maximizer-contraction} gives \(S_h\preceq I\).
For a real unit top eigenvector \(z\) of \(\Sigma\),
\[
 1\ge z^TS_hz
 =t\lambda_1(\Sigma)-t\int p_h(v)(v\cdot z)^2\,d\mu
 \ge t\lambda_1(\Sigma)-tL^2m_h.
\]
Consequently
\(m_h\ge(t\lambda_1(\Sigma)-1)/(tL^2)>0\), uniformly over the
maximizing set. Lemma~\ref{sp:gap}, Theorem~\ref{mf:criterion}, and
the endpoint identity \eqref{mf:ratio} now give
\eqref{sp:positive-margin}.

For arbitrary changing orders, every subsequence has a further
subsequence along which the ordered empirical measures converge
weakly on the compact set \([0,1]\times\overline B(0,L)\).
The limit has Lebesgue position marginal and covariance \(\Sigma\).
The lower bound just proved is the same for every such limit, because
it depends only on \(t,L,r,\lambda_1(\Sigma)\). A subsequence
violating \eqref{sp:positive-margin} would contradict its own
ordered-limit bound.
\end{proof}

\begin{corollary}[A necessary three-eigenvalue condition]
\label{sp:negative-necessity}
Under the bounded-row and covariance-convergence assumptions, if
\[
 \liminf_{N\to\infty}
 \frac{D_1(A_N)}{N(N-1)\per A_N}<0,
\]
then \(t\lambda_3(\Sigma)>1\). Thus a negative leading normalized
endpoint derivative requires at least three limiting nontrivial
eigenvalues of \(A_N\) strictly above two.
\end{corollary}
\begin{proof}
If \(t\lambda_1(\Sigma)>1\), this is the contrapositive of
Theorem~\ref{sp:main}. If \(t\lambda_1(\Sigma)\le1\), then
\(\log(1+s)\le s\), with strict inequality for \(s>0\), shows
\(\Phi(h)<0\) for every \(h\ne0\). If the row projection vanishes
almost everywhere, the same conclusion follows from
\(\Phi(h)=-\|h\|^2\). Thus the maximizing set is \(\{0\}\) and
\(\mathcal R(0)=1\). Theorem~\ref{mf:criterion} gives convergence
of the displayed normalized derivative to zero, including the first
critical boundary \(t\lambda_1(\Sigma)=1\). Compactness of ordered
empirical subsequences gives the same conclusion for changing orders.
\end{proof}

\begin{corollary}[Latent rank at most two]\label{sp:rank-two}
For real rows in \(\mathbb R^2\), if \(t\lambda_1(\Sigma)>1\),
then every sequence of orders satisfies
\[
 \liminf_{N\to\infty}
 \frac{D_1(A_N)}{N(N-1)\per A_N}
 \ge\frac{\bigl(t\lambda_1(\Sigma)-1\bigr)^2}{128t^2L^4}>0.
\]
Together with Theorem~\ref{mf:subcritical} and
Corollary~\ref{mf:arbitrary-orders}, this proves eventual endpoint
increase in both noncritical regimes when \(\Sigma\ne0\).
\end{corollary}

\subsection{Scope and the rank-two literature}\label{sp:prior}

The first critical boundary can still leave derivatives of smaller
order than \(N^2\per A_N\); the sign of such derivatives is not
decided here. A public manuscript of liqsweep
\cite{Liqsweep2026} claims permanent-on-top for all real PSD matrices
of rank at most two. If that claim is accepted, its hereditary form,
closure under nonnegative diagonal addition, and the endpoint
sign-Gram identity already imply finite-order endpoint nonnegativity
for that class plus a nonnegative diagonal. The proofs above use no
result from that manuscript. They give an explicit strict asymptotic
margin and apply in arbitrary latent dimension through the number of
supercritical covariance eigenvalues.
